% Standalone XeLaTeX source; standard commands and distributed packages only.
% Source: public/notes/物工/問題/2019ap_exam.pdf, PDF pp. 2--3.
\documentclass[11pt,a4paper]{article}
\usepackage[margin=24mm]{geometry}
\usepackage{fontspec}
\usepackage{xeCJK}
\setCJKmainfont{HaranoAjiMincho-Regular.otf}
\setCJKsansfont{HaranoAjiGothic-Medium.otf}
\usepackage{amsmath,amssymb}
\usepackage{tikz}
\usetikzlibrary{arrows.meta}
\setlength{\parindent}{1em}
\setlength{\parskip}{0.35em}
\renewcommand{\figurename}{図}
\title{物理工学専攻 院試解答集 ver2\\2019年 物理学 第1問}
\author{政岡凜太郎}
\date{}
\begin{document}
\maketitle

\section*{[1] 固定端をもつ質点ばね系の固有振動}
質点 $n$ の変位を $x_n(t)$ とし、固定端の条件を便宜上
$x_0(t)=x_{N+1}(t)=0$ と書く。質点 $n$ に左右のばねが及ぼす力は
$-k(x_n-x_{n-1})$ と $-k(x_n-x_{n+1})$ なので、
\begin{equation}
  m\frac{\mathrm{d}^2x_n}{\mathrm{d}t^2}
  =-k(2x_n-x_{n-1}-x_{n+1}),\qquad n=1,\ldots,N.
  \label{eq:motion-components}
\end{equation}
従って、$\boldsymbol{x}=(x_1,\ldots,x_N)^\mathsf{T}$ に対する行列は
\begin{equation}
 K=k\begin{pmatrix}
 2&-1&0&\cdots&0\\
 -1&2&-1&\ddots&\vdots\\
 0&-1&2&\ddots&0\\
 \vdots&\ddots&\ddots&\ddots&-1\\
 0&\cdots&0&-1&2
 \end{pmatrix},\qquad
 m\frac{\mathrm{d}^2\boldsymbol{x}}{\mathrm{d}t^2}=-K\boldsymbol{x}.
 \label{eq:K}
\end{equation}
提案された固有モードの内点成分について、恒等式
$\sin(q(n-1))+\sin(q(n+1))=2\sin(qn)\cos q$ より
\begin{align}
 (K\boldsymbol{u}_\ell)_n
 &=k\bigl(2(u_\ell)_n-(u_\ell)_{n-1}-(u_\ell)_{n+1}\bigr)\\
 &=2k(1-\cos q_\ell)(u_\ell)_n.
\end{align}
$n=1,N$ でも、境界値 $(u_\ell)_0=(u_\ell)_{N+1}=0$ が
$\sin 0=0$ および $\sin(q_\ell(N+1))=\sin(\pi\ell)=0$ に
対応するため同じ式が成り立つ。したがって
\begin{equation}
  K\boldsymbol{u}_\ell=\lambda_\ell\boldsymbol{u}_\ell,
  \qquad \lambda_\ell=2k(1-\cos q_\ell),
  \qquad q_\ell=\frac{\pi\ell}{N+1}.
  \label{eq:eigenvalue}
\end{equation}

モード $\ell$ を規格化するには、$\ell$ を固定して $N$ 個の質点成分の
二乗和を求める。$\sin^2 z=(1-\cos 2z)/2$ を用い、問題文の和では
$n=N+1$ の項が $1$ であることに注意すると
\begin{align}
 \sum_{n=1}^{N}\sin^2(q_\ell n)
 &=\frac{N}{2}-\frac12\sum_{n=1}^{N}\cos(2q_\ell n)\\
 &=\frac{N}{2}-\frac12(-1)=\frac{N+1}{2}.
\end{align}
規格化定数を正に選ぶと
\begin{equation}
  (u_\ell)_n=N_\ell\sin(q_\ell n),\qquad
  N_\ell=\sqrt{\frac{2}{N+1}},\qquad
  \sum_{n=1}^{N}(u_\ell)_n^2=1.
  \label{eq:normalization}
\end{equation}
式 \eqref{eq:eigenvalue} の固有値は $0<q_\ell<\pi$ で互いに異なる。
$K$ は実対称行列なので、異なる固有値に属する固有ベクトルは直交する。
$N$ 次元空間にこのような規格化固有ベクトルが $N$ 個あるため、完全
正規直交基底をなす。

$\boldsymbol{x}(t)=\sum_{\ell=1}^{N}\alpha_\ell(t)\boldsymbol{u}_\ell$ と展開し、
固有基底上でベクトル運動方程式を成分ごとに比較すると
\begin{equation}
  \frac{\mathrm{d}^2\alpha_\ell}{\mathrm{d}t^2}
  =-\frac{\lambda_\ell}{m}\alpha_\ell
  =-\omega_\ell^2\alpha_\ell,
  \qquad
  \omega_\ell=\sqrt{\frac{2k}{m}(1-\cos q_\ell)}
  =2\sqrt{\frac{k}{m}}\sin\frac{q_\ell}{2}.
  \label{eq:omega}
\end{equation}
最後の形では $0<q_\ell/2<\pi/2$ なので正の正弦を用いた。分散曲線は
$0<q<\pi$ で単調増加し、二階微分が負なので上に凸である。小さい
$q$ では $\omega\simeq\sqrt{k/m}\,q$ となり、$q=\pi$ での極限値は
$2\sqrt{k/m}$ である。許されるモードは、この曲線上の
$q_\ell=\pi\ell/(N+1)$ に対応する $N$ 個の点である。

\begin{figure}[htbp]
\centering
\begin{tikzpicture}[x=1.45cm,y=0.95cm,>=Stealth]
  \draw[->] (-0.12,0)--(3.45,0) node[right] {$q$};
  \draw[->] (0,-0.12)--(0,2.45) node[above] {$\omega/\sqrt{k/m}$};
  \draw[thick,blue,domain=0:180,samples=100,variable=\t]
    plot ({\t*pi/180},{2*sin(\t/2)});
  \foreach \q in {30,60,90,120,150} {
    \fill[red] ({\q*pi/180},{2*sin(\q/2)}) circle (1.4pt);
  }
  \draw (0,0)--(-0.04,0) node[left] {$0$};
  \draw (1.5708,0.06)--(1.5708,-0.06) node[below=2pt] {$\pi/2$};
  \draw (3.1416,0.06)--(3.1416,-0.06) node[below=2pt] {$\pi$};
  \draw (-0.06,2)--(0.06,2) node[left=2pt] {$2$};
  \node[blue] at (2.25,1.35) {$2\sin(q/2)$};
  \node[red] at (2.48,0.34) {許容モードの例};
\end{tikzpicture}
\caption{固定端鎖の分散関係。赤点は $N=5$ の許容波数を示す例である。}
\label{fig:dispersion}
\end{figure}

\section*{[2] 外力に対する静的応答}
釣り合いでは慣性項が消え、$K\boldsymbol{x}=\boldsymbol{f}$ となる。固有ベクトル
で両辺を展開すると
\begin{equation}
  \sum_{\ell=1}^{N}\lambda_\ell\alpha_\ell\boldsymbol{u}_\ell
  =\sum_{\ell=1}^{N}\beta_\ell\boldsymbol{u}_\ell,
  \qquad
  \lambda_\ell\alpha_\ell=\beta_\ell.
\end{equation}
従って $\alpha_\ell=\beta_\ell/\lambda_\ell$ であり、各モードの変位は外力係数に比例する。その比例係数 $1/\lambda_\ell$ は外力の大きさによらない。各モードの外力係数が $\beta_\ell\ne0$ のとき、設問の応答比はこの比例係数に等しく
\begin{equation}
  \frac{\alpha_\ell}{\beta_\ell}
  =\frac{1}{\lambda_\ell}
  =\frac{1}{m\omega_\ell^2}.
  \label{eq:response}
\end{equation}
外力係数がゼロのモードでは変位係数もゼロとなるが、比例係数 $1/\lambda_\ell$ は同じである。以下では、この外力に依存しない比例係数をモード間で比較する。
$q_\ell$ は $\ell$ とともに増加し、$1-\cos q_\ell$ も狭義単調増加する。
よって $\lambda_\ell$ が最小となるのは $\ell=1$ だけであり、応答比の
最大値は
\begin{align}
 \max_{1\le\ell\le N}\frac{1}{\lambda_\ell}
 &=\frac{1}{4k\sin^2\!\left(\dfrac{\pi}{2(N+1)}\right)}\\
 &\simeq\frac{(N+1)^2}{\pi^2k}\qquad(N\gg1),
 \label{eq:max-response}
\end{align}
である。近似では $\sin z\simeq z$ を用いた。

\section*{[3] 各質点に加えた調和ポテンシャル}
追加ポテンシャル $\tfrac12 k_0x_n^2$ は質点 $n$ に復元力 $-k_0x_n$ を
加える。したがって運動方程式の剛性行列は
\begin{equation}
  K_{0}=K+k_0I.
\end{equation}
単位行列は全てのベクトルを同じ向きのまま保つため、固有ベクトルは
$\boldsymbol{u}_\ell$ のままで、固有値だけが $k_0$ 増える。
\begin{equation}
  \lambda_{\ell,0}=2k(1-\cos q_\ell)+k_0
  =4k\sin^2\frac{q_\ell}{2}+k_0,
  \qquad
  \omega_\ell=\sqrt{\frac{4k\sin^2(q_\ell/2)+k_0}{m}}.
  \label{eq:pinned-omega}
\end{equation}
この式は $k_0$ と $k$ が同じ単位（力/長さ）であることとも整合する。
有限の $N$ では最小値は $\ell=1$ であり、$q_1=\pi/(N+1)\to0$ より
\begin{equation}
  \lim_{N\to\infty}\min_{1\le\ell\le N}\omega_\ell
  =\sqrt{\frac{k_0}{m}}.
\end{equation}

図\ref{fig:pinned-dispersion}は例として $k_0=k$ とした曲線を描く。
一般には縦軸を $\sqrt{k/m}$ で規格化した式
$\omega/\sqrt{k/m}=\sqrt{4\sin^2(q/2)+k_0/k}$ に従い、$q=0$ での
縦軸切片は $\sqrt{k_0/k}$ となる。従って正の $k_0$ があると、長波長
極限にも有限のギャップが残る。

\begin{figure}[htbp]
\centering
\begin{tikzpicture}[x=1.45cm,y=0.95cm,>=Stealth]
  \draw[->] (-0.12,0)--(3.45,0) node[right] {$q$};
  \draw[->] (0,-0.12)--(0,2.75) node[above] {$\omega/\sqrt{k/m}$};
  \draw[thick,blue,domain=0:180,samples=100,variable=\t]
    plot ({\t*pi/180},{sqrt(1+4*sin(\t/2)^2)});
  \foreach \q in {30,60,90,120,150} {
    \fill[red] ({\q*pi/180},{sqrt(1+4*sin(\q/2)^2)}) circle (1.4pt);
  }
  \draw (1.5708,0.06)--(1.5708,-0.06) node[below=2pt] {$\pi/2$};
  \draw (3.1416,0.06)--(3.1416,-0.06) node[below=2pt] {$\pi$};
  \draw (-0.06,1)--(0.06,1) node[left=2pt] {$1$};
  \draw (-0.06,{sqrt(5)})--(0.06,{sqrt(5)}) node[left=2pt] {$\sqrt{5}$};
  \node[blue] at (2.13,2.15) {$k_0=k$};
  \node[red] at (2.57,0.7) {許容モードの例};
\end{tikzpicture}
\caption{$k_0=k$ の場合の分散関係。赤点は $N=5$ の許容波数を示す例である。}
\label{fig:pinned-dispersion}
\end{figure}
\end{document}
